\section{Introduction}
\label{sec:introduction}

Two models can be close to one another and still play different roles in a larger collection. One may be reconstructed almost as well by a closely related peer as by all available peers, while reconstructing that related model benefits from additional models. Conversely, a target can differ from every individual peer yet be closely approximated by their combination. The same pairwise comparison cannot distinguish dependence on one peer from dependence on the group. This distinction concerns \emph{who covers whom}: how the collection reproduces a target's responses, and which peers make that reproduction possible.

A reconstruction score alone also leaves this structure unresolved. Two targets with similar residuals may rely on very different subsets of the ecosystem. Removing a single related checkpoint can disrupt the approximation of one target while scarcely affecting the other. The relevant comparison therefore combines reconstruction error with the change in error when the available peers change. It gives an auditor a target-specific account of where the collection's coverage comes from.

We formalize \emph{behavioral coverage} through recorded responses on matched inputs. An audit specifies the response to observe, the peer set, and the fitting and evaluation conditions. A fixed convex combination uses the same non-negative, sum-to-one weights on every audited input. This asks whether the peers' existing response traces jointly reproduce the target on their original scale. The Peer-Inexpressible Residual (\PIER) is the remaining target--peer difference. \DISCO fits the combination on one sample and evaluates residual magnitude on another. Group gain compares that fit with a single peer selected on the fitting sample; peer-removal effects measure the loss of coverage after each peer is excluded and the combination refitted.

Our principal experiment studies eight instruction- and reasoning-tuned language models on MMLU-Pro \citep{wang2024mmlupro}. The primary response is each model's reference-answer probability, balanced over cyclic option rotations. One weight vector must reconstruct its variation across questions and stress conditions. Six targets benefit from a peer group under both stress families and both matched fitting losses. Qwen2.5 has a different structure: General Reasoner, its reasoning-tuned sibling, achieves essentially the full group's approximation quality. In the reverse direction, the Qwen sibling remains important, but additional peers reduce General Reasoner's error by $15$--$18\%$. Mistral exhibits a $19$--$20\%$ group advantage. The weak removal effects within the Phi pair provide a contrast: lineage identifies candidate relations, while measured responses determine their strength.

We examine this structure through a second response view that retains the probabilities of \emph{all} answer options. In the canonical-order vector analysis, the Qwen pair again differs in whether additional peers improve reconstruction, and Mistral retains a collective advantage. Gemma, however, has different group-gain signs in the two settings. The response map thus selects a substantive object of comparison: agreement in reference-answer confidence and agreement across the complete option distribution can reveal different coverage relations.

Matched stress adds a further dimension. Irrelevant context increases the balanced residual for Qwen2.5 and phi-4 while decreasing it for Mistral and OLMo. The vision experiments connect such changes to the supporting peer combination and to a small set of shared high-residual images. Forecasting data illustrate the separate role of task outcomes: a response difference can accompany either an improvement or a deterioration in prediction error. Together, these analyses distinguish the amount, source, and location of a collection's coverage.

The estimation step builds on convex response fitting, as in synthetic control \citep{abadie2010synthetic}. Here the observations are matched model responses, and the object of study is their reconstruction relation. This differs from explanations of a single model \citep{ribeiro2016lime,sundararajan2017integrated,wachter2017counterfactual}, comparisons of internal representations \citep{raghu2017svcca,kornblith2019cka}, and combinations optimized for task prediction \citep{hinton2015distilling,wortsman2022model}.

\noindent\textbf{Contributions.}
\begin{itemize}[leftmargin=1.2em,itemsep=2pt,topsep=3pt]
\item \emph{Collective response coverage.} We formulate target reconstruction by a fixed peer-response class and combine held-out \PIER, group gain, and removal effects to distinguish the amount of coverage from its source. Projection properties support the estimator even when coefficient representations are non-unique.
\item \emph{Directional and distributed relations.} An eight-model language-model audit separates single-sibling coverage from sibling-supported and distributed approximation. A complete-option-distribution diagnostic examines which of these patterns extend beyond reference-answer confidence.
\item \emph{Stress and localization.} Matched interventions move target residuals in opposite directions within one ecosystem. Response balancing tests those directions, and vision audits identify changes in peer support and concentrated shared residuals.
\end{itemize}
